\documentclass[12pt,a4paper]{article}

\usepackage[T2A]{fontenc}
\usepackage[utf8]{inputenc}
\usepackage[russian]{babel}
\usepackage{amsmath,amssymb}
\usepackage[a4paper,margin=2cm]{geometry}

\setlength{\parindent}{0pt}
\setlength{\parskip}{0.45em}

\begin{document}

\begin{center}
{\large\bfseries Экзаменационная программа}

по дисциплине «Многомерный анализ, интегралы и ряды»

весенний семестр 2025--2026 учебного года

для потока Сакбаева В.Ж.
\end{center}

\begin{enumerate}
\item Предел последовательности точек в $\mathbb{R}^n$. Связь между сходимостью
последовательности точек и сходимостью последовательностей их координат.
Внутренние, предельные, изолированные точки множества. Открытые и замкнутые
множества, их свойства. Непрерывная кривая, область, компакт. Критерий
компактности в $\mathbb{R}^n$. Внутренность, замыкание и граница множества.

\item Предел числовой функции нескольких переменных. Предел функции по множеству.
Пределы по направлениям. Непрерывность функции нескольких переменных в точке и
по множеству. Непрерывность сложной функции. Свойства функций, непрерывных на
компакте: ограниченность, достижимость точных нижней и верхней граней,
равномерная непрерывность. Теорема о промежуточных значениях функции,
непрерывной в области.

\item Частные производные функции нескольких переменных. Дифференцируемость
функции в точке, дифференциал. Необходимые условия дифференцируемости,
достаточные условия дифференцируемости функции нескольких переменных.
Дифференцируемость сложной функции. Инвариантность формы дифференциала
относительно замены переменных. Градиент и производная по направлению.

\item Частные производные высших порядков. Независимость смешанной частной
производной от порядка дифференцирования. Дифференциалы высших порядков.
Формула Тейлора для функций нескольких переменных с остаточным членом в формах
Лагранжа и Пеано.

\item Определение измеримости по Жордану множества в $n$-мерном евклидовом
пространстве. Критерий измеримости. Измеримость объединения, пересечения и
разности измеримых множеств. Конечная аддитивность меры Жордана.

\item Определенный интеграл Римана. Верхние и нижние суммы Дарбу. Свойства сумм
Дарбу. Критерии интегрируемости. Интегрируемость непрерывной функции,
монотонной функции, ограниченной функции с конечным числом точек разрыва.
Аддитивность интеграла по отрезкам, линейность интеграла, интегрируемость
произведения функций, интегрируемость модуля интегрируемой функции,
интегрирование неравенств, теорема о среднем. Свойства интеграла с переменным
верхним пределом: непрерывность, дифференцируемость. Формула
Ньютона--Лейбница. Замена переменных и интегрирование по частям в определенном
интеграле.

\setcounter{enumi}{7}
\item Геометрические приложения определенного интеграла: площадь криволинейной
трапеции, объем тела вращения, длина кривой. Площадь поверхности вращения (без
доказательства).

\item Несобственный интеграл. Критерий Коши сходимости интеграла. Интегралы от
знакопостоянных функций, признак сравнения. Интегралы от знакопеременных
функций, сходимость и абсолютная сходимость. Признаки Дирихле и Абеля
сходимости интегралов.

\item Числовые ряды. Критерий Коши сходимости ряда. Знакопостоянные ряды:
признак сравнения, признаки Даламбера и Коши, интегральный признак.
Знакопеременные ряды, сходимость и абсолютная сходимость. Признаки Лейбница,
Дирихле и Абеля. Независимость суммы абсолютно сходящегося ряда от порядка
слагаемых. Произведение абсолютно сходящихся рядов.

\item Равномерная сходимость функциональных последовательностей и рядов.
Критерий Коши равномерной сходимости. Непрерывность суммы равномерно
сходящегося ряда из непрерывных функций. Интегрирование и дифференцирование
функциональных последовательностей и рядов. Признаки Вейерштрасса, Дирихле и
Абеля равномерной сходимости функциональных рядов.

\item Степенные ряды с комплексными членами. Круг и радиус сходимости. Характер
сходимости степенного ряда в круге сходимости. Формула Коши--Адамара.
Сохранение радиуса сходимости степенного ряда при почленном дифференцировании
и интегрировании ряда. Первая теорема Абеля.

\item Степенные ряды с действительными членами. Почленное интегрирование
степенного ряда. Бесконечная дифференцируемость суммы степенного ряда на
интервале сходимости. Единственность представления функции степенным рядом.
Достаточные условия разложимости бесконечно дифференцируемой функции в
степенной ряд. Ряд Тейлора. Формула Тейлора с остаточным членом в интегральной
форме. Пример бесконечно дифференцируемой функции, не разлагающейся в
степенной ряд. Разложение в ряд Тейлора основных элементарных функций:
$e^x$, $\cos x$, $\sin x$, $\ln(1+x)$, $(1+x)^\alpha$. Разложение в степенной
ряд функции $e^z$ комплексного переменного $z$.

\item Экстремумы функций многих переменных: необходимое условие, достаточное
условие.
\end{enumerate}

\end{document}
